\begin{itemize}
    \item 1- Your jitter curves here are different than what you have presented in the final presentation or other slides you have sent me. What changed? 
    \item 1A- In the presentation, I showed a $425 ps$ of FWHM corresponding to the pulsed laser operating at $1MHz$ and $33 \mu W$. On the report opted to show a $160 ps$ FWHM at $80MHz$ and $33 \mu W$ . I saved all the data in: \verb|datacoincidence_meas_run1_80MHz_100uW| and \verb|data_coincidence_meas_run2_1MHz_33uW|. I also have saved values for if you are intrested in exporing these results. I also have: \verb|data_coincidence_meas_run1_1MHz_3uW| and \verb|data_coincidence_meas_run1_ThermalSources_g2|. I placed the two figures bellow for referece. In my opinion, this is due to the source jitter rather than the SPAD. 
    \item 2- These look better, but how did you get these, whats the difference? 
    \item 2A- Both the power and the frequency changed. Although I suspect the frequency is the main culpit. Longer jitter corresponded to the setting with lower frequency. 
    \item 3- Important for full Coincidence Measurements section!: I believe you did the same mistake
here as during your final presentation. You should have used two SPADs in this
experiment, as written in the procedures I gave you in the beginning. Therefore,
there is no auto-correlation, there is no term of auto-correlation function for
laser or thermal lamp etc. here. It is still cross-correlation. You should have
built a histogram based on timestamps between two SPAD channels, not with
the laser. As I wrote in the presentation, there is no laser sync used here. Or, I
am very confused with the experiment you performed here. You switched to
getting auto-correlation of single SPADs (not laser) with rotating diffuser
experiment. You utilized SPADA and SPADB until that experiment, which makes
your results cross-correlation until rotating diffuser. 
\item 3A-You are entirely correct —the measurements were \textbf{cross-correlations} between the two SPADs, not autocorrelations. I had the Wiener-Khinchin theorem on my mind from my recent metrology exam and mixed up the terminology in the text. I proceed to clarify this mistake with the actual syntax I used for the measurements:
\begin{itemize} \footnotesize
    \item For the 1st Measurement, I measured cross-correlation: 
    \item[] $counter = tt.Counter(tagger, [spadA, spadB, sync], binwidth=exposuretime, nvalues=1) $
    \item[] $coinc = tt.Coincidence(tagger, [spadA, spadB], coincidencewindow) $
    \item[] $coinccounter = tt.Counter(tagger, [coinc.getChannel()], binwidth=exposuretime, nvalues=1)$
    \item For the 2nd Measurement and the jitter curves, again is cross-correlation: 
    \item[] $tt.Correlation(tagger, spadA, spadB, binwidth=binwidthps, nbins=nbins)$
\end{itemize}

\end{itemize}


\begin{center}
    \includegraphics[width=0.4\textwidth]{Figures/answers_jitter1MHZ.png}
    \includegraphics[width=0.4\textwidth]{Figures/answers_jitter80MHZ.png}
\end{center}


